\section{Introduction}
\label{sec:introduction}
Graphic design involves the purposeful organization of typography, imagery, vector shapes, colors, and spatial layout to communicate information and visual intent \citep{huang2023survey,zou2025fragment}. Unlike unconstrained natural-image synthesis, graphic-design applications often require coordinated text and non-text components, legibility, visual hierarchy, and design constraints.

\paragraph{Application tasks.}
Computational research addresses several concrete, but non-equivalent, design tasks: \emph{poster layout and composition} \citep{posterlayout2023,vinci2021,posteragent2026}; \emph{advertising and marketing layout generation} \citep{cglgan2022,layoutdetr2024,autoposter2023}; \emph{presentation-slide creation and refinement} \citep{autopresent2025,pptagent2025}; \emph{infographic generation and editing} \citep{bizgen2025,yang2026infoedit}; and \emph{scientific illustration and paper-to-poster design} \citep{autofigure2026,postergen2026}. Specialized research also studies \emph{typography and vector icons} \citep{textdiffuser2023,wordasimage2023,iconshop2023}. These citations support specific studied tasks rather than treating all examples as interchangeable instances of a single image-generation problem.

\paragraph{Artifact and workflow.}
A poster-layout system may return the bounding boxes of text and logo elements, a slide-generation system may output an editable program or document, and an infographic-editing system may modify an existing visual composition \citep{posterlayout2023,autopresent2025,yang2026infoedit}. Consequently, the application category alone does not determine editability. For example, a designer may need to replace a headline, resize the canvas, move a logo, or export independently manipulable components. A flattened bitmap, an SVG path sequence, and an HTML document can render similarly yet expose different post-generation operations.

These applications have prompted multiple technological approaches. LayoutDM and LayoutFlow model element geometry \citep{layoutdm2023,layoutflow2024}; content-aware systems use image context for composition \citep{cglgan2022,layoutdetr2024}. Hierarchical design pipelines separate intention, planning, and rendering \citep{cole2023,opencole2024}. Agentic and code-native systems generate structured artifacts that can undergo further changes \citep{designascode2026,editablevisual2026}, while specialized typography models target readable text and semantically meaningful lettering \citep{textdiffuser2023,wordasimage2023}. Vector synthesis provides another form of structural control \citep{vectorfusion2023,iconshop2023}.

These streams are difficult to compare by model architecture alone. A diffusion model may output a bitmap, geometry, or vector paths; an agent may return a rendered image or a structured design file. Two studies can both claim editability even if one demonstrates manual SVG path manipulation and another measures a human's efficiency when resizing a poster. Moreover, high-quality rendered output does not by itself imply that an edit instruction can change only its intended target.

Prior surveys discuss graphic design intelligence and the wider AI-driven graphic design pipeline \citep{huang2023survey,zou2025fragment}. This survey develops a complementary analytical perspective through three questions:
\begin{enumerate}[leftmargin=*,nosep]
\item \textbf{Representation (R):} What artifact, geometry, or program is delivered?
\item \textbf{Editability (E):} What operation can be applied to which design object?
\item \textbf{Evaluation (V):} What outcome is measured, using what protocol and evidence?
\end{enumerate}
The proposed R--E--V crosswalk connects representations to supported operations and tests, rather than assigning an undifferentiated editable/not-editable label.

Our contributions are an orthogonal coding framework, an evidence-sensitive synthesis of representative technical lines, and a concrete research agenda for testing edit instructions, non-target preservation, output validity, and design workflow utility. We explicitly distinguish primary-paper evidence from abstract-level claims. Figures~\ref{fig:rev} and \ref{fig:protocol} present the conceptual distinctions, while Sec.~\ref{sec:related} compares the organization and visualization patterns of existing surveys. This first manuscript is a structured narrative synthesis; completeness and screening reproducibility must be established in later revisions.
